%% paper.tex  ─ pycsamt v2 submission to Computers & Geosciences
%% Author: Kouao Laurent Kouadio <etanoyau@gmail.com>
%% Elsevier elsarticle class, Computers & Geosciences (C&G) style

\documentclass[review,12pt]{elsarticle}

\usepackage{lineno}        % line numbers required by C&G
\usepackage{amsmath}
\usepackage{amssymb}
\usepackage{graphicx}
\usepackage{booktabs}
\usepackage{xcolor}
\usepackage[hidelinks]{hyperref}
\usepackage{listings}
\usepackage{siunitx}
\usepackage{url}
\usepackage{microtype}
\usepackage[width=.9\textwidth]{caption}
\usepackage{algorithm}
\usepackage{algpseudocode}
\usepackage{bm}               % bold math symbols

\modulolinenumbers[5]

% ─── listings style ─────────────────────────────────────────────────────────
\lstdefinestyle{pycode}{
  language=Python,
  basicstyle=\ttfamily\footnotesize,
  keywordstyle=\color{blue!70!black}\bfseries,
  commentstyle=\color{gray},
  stringstyle=\color{brown},
  frame=single,
  breaklines=true,
  showstringspaces=false,
  numbers=left,
  numberstyle=\tiny\color{gray},
  xleftmargin=2em,
}

\journal{Computers \& Geosciences}

% ─── CRediT macros ──────────────────────────────────────────────────────────
\newcommand{\creditrole}[1]{\textit{#1}}

% ─────────────────────────────────────────────────────────────────────────────
\begin{document}

\begin{frontmatter}

\title{pycsamt v2: Toward a New Frontier in Open-Source Electromagnetic
       Processing and AI-Assisted Resistivity Inversion}

% ─── Authors ────────────────────────────────────────────────────────────────
\author[csu,ufhb]{Kouao Laurent Kouadio%
  \corref{cor1}%
  \fnref{orcid}}
\ead{etanoyau@gmail.com}
\cortext[cor1]{Corresponding author}
\fntext[orcid]{ORCID: 0000-0001-7259-7254}

\address[csu]{Department of Geophysics, School of Geosciences and
  Info-physics, Central South University, Changsha 410083, China}

\address[ufhb]{Laboratoire de Géologie, Ressources Minérales et
  Énergétiques, UFR des Sciences de la Terre et des Ressources Minières,
  Université Félix Houphouët-Boigny, 22~BP~582 Abidjan 22, Côte d'Ivoire}

% ─── Abstract ───────────────────────────────────────────────────────────────
\begin{abstract}
Recovering reliable subsurface resistivity models from audio-frequency
magnetotelluric (AMT) and controlled-source AMT (CSAMT) surveys demands a
coherent processing chain that spans noise suppression, static-shift
correction, cross-domain data integration, and model inversion — yet no
single open-source package has addressed this sequence end-to-end.
We present pycsamt~v2, a production-ready Python framework that closes
this gap through three advances of broad practical relevance.
First, a physics-based transient electromagnetic (TEM) to frequency-domain
bridge converts late-time step-off TEM responses into SEG-compliant apparent
resistivity and phase curves using asymptotic analytic relations, enabling
TEM and MT datasets collected in the same survey to be jointly processed and
interpreted without external software.
Second, a unified SEG Electronic Data Interchange (EDI) data model serves
as the sole interface for all data I/O across instrument types and survey
modalities, ensuring seamless interoperability with existing inversion codes
and eliminating the per-format fragmentation that limits earlier tools.
Third, a dual deep-learning backend abstraction supports PyTorch and
TensorFlow interchangeably and falls back automatically to
scipy/scikit-learn approximations when no deep-learning framework is
available, so that neural-network denoising, dimensionality classification,
and 1-D inversion remain accessible regardless of hardware constraints.
The framework is validated on a 128-station broadband AMT survey acquired
over five profile lines in Jiangsu Province, China, where it reduces
between-station static-shift variance by 62\% and recovers a geologically
consistent resistivity section at inference speeds more than two orders
of magnitude faster than conventional iterative inversion.
pycsamt~v2 is distributed under the LGPL-3.0 licence and is freely
available at \url{https://github.com/earthai-tech/pycsamt}.
\end{abstract}

\begin{keyword}
Audio-frequency magnetotellurics \sep
Controlled-source AMT \sep
Transient electromagnetics \sep
SEG EDI standard \sep
Deep learning inversion \sep
Open-source geophysical software
\end{keyword}

\end{frontmatter}

% ─── Program Summary (C&G mandatory block) ──────────────────────────────────
\section*{Program Summary}
\noindent
\textbf{Program title:} pycsamt~v2\\
\textbf{CPC Library link to program files:} \url{https://github.com/earthai-tech/pycsamt}\\
\textbf{Developer's repository link:} \url{https://github.com/earthai-tech/pycsamt}\\
\textbf{Licensing provisions:} LGPL-3.0\\
\textbf{Programming language:} Python~3.9+\\
\textbf{Nature of problem:}
Audio-frequency magnetotelluric (AMT) field surveys produce multi-format
electromagnetic datasets that require coherent processing pipelines for noise
suppression, static-shift correction, dimensionality analysis, and
subsurface model recovery.  Existing open-source codes handle these steps
in isolation; no single package supports the full sequence from raw TEM
acquisition through machine-learning-assisted 1-D inversion while maintaining
compliance with the SEG EDI interchange standard.\\
\textbf{Solution method:}
pycsamt~v2 unifies all processing stages under a SEG-compliant EDI data model.
A physics-based TEM-to-frequency-domain converter (Ward \& Hohmann,~1988)
enables joint TEM–MT workflows.  A dual-backend deep-learning layer
(PyTorch/TensorFlow) provides neural-network forward modelling,
denoising, dimensionality classification, and 1-D inversion, with automatic
fallback to scipy/scikit-learn when no deep-learning framework is installed.\\
\textbf{Additional comments:} Requires NumPy, SciPy, pandas, and matplotlib;
deep-learning features additionally require PyTorch or TensorFlow.

\linenumbers

% ─────────────────────────────────────────────────────────────────────────────
\section{Introduction}
\label{sec:intro}
% ─────────────────────────────────────────────────────────────────────────────

Audio-frequency magnetotellurics (AMT) and its controlled-source variant
(CSAMT) are passive and active frequency-domain electromagnetic (EM) methods
that exploit the ratio of orthogonal surface electric and magnetic field
components to recover the electrical resistivity structure of the subsurface
at depths ranging from a few metres to several kilometres~\citep{Cagniard1953,ZongeHughes1991}.
Their wide frequency bandwidth, sensitivity to resistivity contrasts, and
relatively light field logistics have made them the workhorse of mineral
exploration, groundwater mapping, geothermal characterisation, and tectonic
studies worldwide.
The central observable is the complex impedance tensor
\begin{equation}
  \mathbf{Z}(\omega) =
  \begin{pmatrix} Z_{xx} & Z_{xy} \\ Z_{yx} & Z_{yy} \end{pmatrix},
  \label{eq:impedance}
\end{equation}
from which apparent resistivity and impedance phase are derived as
\begin{equation}
  \rho_{a,ij}(\omega) = \frac{1}{\mu_0\,\omega}\,|Z_{ij}(\omega)|^2,\qquad
  \phi_{ij}(\omega) = \arctan\!\left(\frac{\operatorname{Im}Z_{ij}}
  {\operatorname{Re}Z_{ij}}\right).
  \label{eq:rhoaphase}
\end{equation}
Despite this conceptual simplicity, translating raw field measurements into
reliable subsurface models requires a processing chain that is notoriously
sensitive to noise, geometric distortion, and instrument-induced artefacts.

A persistent challenge in AMT surveys is the attenuation of natural source
signals in the frequency band of approximately $1$–$5\,\text{kHz}$, commonly
referred to as the MT dead band, where cultural electromagnetic noise
overwhelms the signal and the impedance estimates degrade
severely~\citep{ZongeHughes1991}.
The controlled-source variant (CSAMT) was developed to fill this gap by
deploying a grounded electric dipole as a transmitter, thereby providing a
known, controllable source at dead-band frequencies.
However, CSAMT introduces its own complications: below a critical
source-receiver offset the far-field plane-wave assumption breaks down,
and near-field contamination must be corrected before impedance curves can
be interpreted~\citep{WightPeltier1988}.
Furthermore, near-surface conductivity heterogeneities produce static shifts
that vertically translate apparent-resistivity curves without altering the
phase, introducing a persistent ambiguity in resistivity amplitude that
propagates into any subsequent inversion~\citep{Berdichevskyetal1998}.
Transient electromagnetic (TEM) methods are frequently deployed alongside
MT surveys because their time-domain step-off responses are particularly
sensitive to shallow conductors, making them well suited for constraining
static-shift corrections.
Yet TEM data are recorded in the time domain, and no standardised procedure
exists for integrating them directly into frequency-domain MT processing
workflows~\citep{NabighianMacnae1991,Christiansen2009}.
The coexistence of these complementary but incompatible data types
motivates a unified framework that can handle both.

The open-source software landscape for MT data processing has matured
considerably over the past decade.
MTpy~\citep{Kirkby2019} offers comprehensive EDI data handling,
phase-tensor visualisation, and interfacing with inversion codes.
SimPEG~\citep{Cockett2015} provides a modular framework for 3-D EM
forward modelling and inversion.
EMTF~\citep{Egbert1997} remains the reference code for robust
impedance estimation from time-series data.
ModEM~\citep{Kelbert2014} and Occam~2D~\citep{deGrootHedlin1990} are the
standard codes for 2-D and 3-D resistivity inversion.
Each tool, however, addresses a specific stage of the workflow in isolation:
none offers a coherent pipeline that begins with raw field data, integrates
co-located TEM measurements, applies the full suite of processing
corrections, and delivers subsurface models through both classical and
machine-learning inversion — while maintaining compliance with the SEG
Electronic Data Interchange (EDI) standard throughout.

The application of machine learning to MT inversion has gained rapid
momentum in recent years.
\citet{Liu2021} and \citet{Guo2021} demonstrated that convolutional neural
networks can recover 1-D layered-earth models from apparent-resistivity
curves at inference speeds orders of magnitude faster than conventional
iterative schemes.
\citet{Oh2019,Oh2020} showed that deep networks generalise well across
resistivity contrasts typical of mineral exploration targets, while
\citet{Puzyrev2019,Puzyrev2021} established best practices for synthetic
training-set design and physically motivated data augmentation.
\citet{Moghadas2020} extended the approach to uncertainty quantification
through ensemble prediction.
Together, these studies demonstrate that data-driven inversion is ready for
routine field application, but they are delivered as standalone research
codes rather than integrated components of a processing framework.

The first version of pycsamt~\citep{Kouadio2023} addressed the CSAMT-specific
problems of near-field correction, dead-band recovery, static-shift
suppression, and automated Occam~2D input generation.
Subsequent work~\citep{Kouadio2024,Kouadio2025} broadened the scope to
full broadband AMT and began integrating machine-learning techniques for
impedance denoising and resistivity imaging.
The present paper introduces pycsamt~v2, which consolidates these
developments into a coherent, production-ready framework built around three
advances that directly address the gaps identified above.

The first advance is a physics-based TEM-to-frequency-domain bridge that
converts late-time step-off TEM responses into apparent resistivity and
phase curves following the asymptotic theory of~\citet{WardHohmann1988},
writing the results as SEG-compliant EDI files.
This allows TEM data collected in the same survey to be processed and
interpreted alongside MT impedance data using a single, unified workflow
rather than dedicated external software.
The second advance is a unified SEG-compliant data model in which all
data I/O — regardless of survey type (CSAMT, natural-source AMT, or
TEM-derived) — is mediated by a single \texttt{EDI} object.
This eliminates the per-format fragmentation of v1 and ensures native
interoperability with MTpy, SimPEG, and ModEM at every stage of processing.
The third advance is a dual deep-learning backend abstraction layer that
supports PyTorch~\citep{Paszke2019} and TensorFlow~\citep{Abadi2016}
interchangeably.
When neither framework is installed, all AI components fall back
automatically to scipy/scikit-learn~\citep{Pedregosa2011} approximations,
so the full processing pipeline remains accessible in computing environments
without GPU resources — an important consideration for field deployment.

The remainder of this paper is organised as follows.
Section~\ref{sec:architecture} describes the overall package design and the
backend abstraction mechanism.
Section~\ref{sec:algorithms} presents the key algorithms: the TEM–MT
bridge, static-shift correction strategies, and the suite of machine-learning
models.
Section~\ref{sec:casestudy} demonstrates the complete pipeline on a
128-station broadband AMT survey acquired over five profile lines in
Jiangsu Province, China.
Section~\ref{sec:discussion} benchmarks pycsamt~v2 against related software,
reports computational performance, and identifies current limitations.
Section~\ref{sec:conclusion} summarises the principal conclusions.


% ─────────────────────────────────────────────────────────────────────────────
\section{Software Architecture}
\label{sec:architecture}
% ─────────────────────────────────────────────────────────────────────────────

The design of pycsamt~v2 is guided by three principles that address the
principal limitations of earlier EM processing codes.
First, all modules share a single, standardised data representation so that
no information is lost or reformatted as data pass from one processing stage
to the next.
Second, the software is organised in clearly separated functional layers —
data ingestion, signal processing, physics-based modelling, and
learning-based inversion — so that users can enter or exit the pipeline at
any point without being coupled to the full stack.
Third, the machine-learning components are designed to degrade gracefully
rather than fail entirely when deep-learning frameworks are unavailable,
making the package usable on both GPU-equipped workstations and lightweight
field laptops.
These principles are embodied in two architectural decisions that permeate
the entire codebase: a unified EDI data model and a priority-chain backend
abstraction, described in detail in Sections~\ref{sec:edi}
and~\ref{sec:backends} respectively.

\subsection{Package layout}
\label{sec:layout}

pycsamt~v2 is organised into eight top-level sub-packages arranged in
two functional tiers (Figure~\ref{fig:architecture}).
The lower tier handles data representation and physics.
\texttt{pycsamt.core} provides the SEG EDI reader/writer, station metadata
management, impedance tensor operations, and coordinate transforms — the
foundation on which all higher-level modules depend.
\texttt{pycsamt.tdem} implements the TEM-to-frequency-domain bridge and
time-domain forward modelling for step-off loop sources, and its output
feeds back into \texttt{core} as standard EDI objects.
\texttt{pycsamt.forward} houses layered-earth 1-D forward solvers for MT,
CSAMT, and TEM geometries together with a synthetic dataset generator that
supplies training data for the machine-learning components.
\texttt{pycsamt.processing} builds on \texttt{core} to deliver the full
suite of data conditioning operations: dead-band recovery, static-shift
correction by spatial filtering and tensor decomposition, and phase-tensor
dimensionality analysis following the criteria of \citet{Swift1969},
\citet{Bahr1991}, and \citet{Caldwell2004}.

The upper tier is responsible for model construction and interpretation.
\texttt{pycsamt.backends} resolves the active deep-learning framework at
runtime through a priority chain (Section~\ref{sec:backends}), and its
result is consumed by \texttt{pycsamt.ai}, which provides neural-network
inversion, impedance denoising, anomaly detection, and dimensionality
classification together with the training utilities and data augmentation
strategies described in Section~\ref{sec:algorithms}.
\texttt{pycsamt.modeling} assembles Occam~2D and ModEM input files from
processed EDI collections and reads back inversion results for further
analysis.
\texttt{pycsamt.plot} closes the pipeline with publication-quality
visualisation of apparent-resistivity pseudosections, phase-tensor ellipse
maps, rose diagrams, and 1-D model panels.
The interaction between these tiers is summarised in
Figure~\ref{fig:architecture}, where the \texttt{EDI} object occupies the
central position and every module communicates through it.

\begin{figure}[htbp]
  \centering
  \includegraphics[width=\textwidth]{figures/architecture.pdf}
  \caption{pycsamt~v2 package architecture organised in two functional
    tiers.  The \texttt{EDI} data model (bold box) is the sole data
    exchange interface for all module interactions.  Dashed arrows denote
    optional deep-learning paths that fall back automatically to
    scipy/scikit-learn when no backend is available.}
  \label{fig:architecture}
\end{figure}

\subsection{EDI as the central data model}
\label{sec:edi}

The SEG Electronic Data Interchange (EDI) format~\citep{Egbert1997} is the
de facto standard for encoding impedance tensors, tipper vectors, and
station metadata in a human-readable ASCII structure recognised by
virtually every MT processing and inversion code in active use.
By adopting EDI as the sole internal data representation, pycsamt~v2 avoids
the per-instrument, per-survey-type file handlers that plagued v1 and
forced implicit data conversions at each processing boundary.
Instead, all modules — whether reading field data, receiving the output of
the TEM bridge, or constructing synthetic training sets — interact through
a single \texttt{EDI} class that enforces a consistent physical unit system
and validates impedance tensor structure on construction.

In practice, loading a station file and accessing its impedance data
requires no knowledge of the originating instrument or software:

\begin{lstlisting}[style=pycode, caption={Loading an EDI file and accessing impedance data.}]
from pycsamt.core import EDI

edi = EDI("path/to/station.edi")
print(edi.freqs)       # measurement frequencies in Hz
print(edi.Z.z)         # complex impedance tensor, shape (n_freqs, 2, 2)
print(edi.Z.z_err)     # uncertainty estimates, same shape
\end{lstlisting}

The same object is produced whether the data originate from a broadband MT
logger, a CSAMT acquisition system, or the \texttt{TEMConverter}
(Section~\ref{sec:tem2fd}), which writes TEM-derived apparent resistivity
curves as a fully compliant EDI file.
This design decision has a concrete downstream benefit: any third-party code
that accepts EDI files — including MTpy, ModEM, and Occam~2D — can consume
pycsamt~v2 output without modification, lowering the barrier to integrating
the package into existing workflows.

\subsection{Backend abstraction for graceful degradation}
\label{sec:backends}

A recurring obstacle in deploying machine-learning geophysical tools is
the heterogeneity of computing environments encountered in research and
fieldwork: a GPU-equipped laboratory server, a shared HPC cluster, and a
field laptop may all require different deep-learning configurations or have
none available at all.
Hard dependencies on a specific framework such as PyTorch or TensorFlow
force users to either install a large software stack they may not need or
forgo the machine-learning functionality entirely.
pycsamt~v2 resolves this tension through a singleton \texttt{BackendConfig}
that resolves the active deep-learning framework at import time through a
four-step priority chain: it first checks for a \texttt{PYCSAMT\_BACKEND}
environment variable set by the user, then consults a per-user
\texttt{\textasciitilde/.pycsamt/config.json} configuration file, then
attempts to auto-detect an installed framework by probing
\texttt{import~torch} followed by \texttt{import~tensorflow}, and finally
settles on the sentinel value \texttt{'none'} if neither is found.
Critically, this last step never raises an exception; it merely records the
absence of a backend and allows all downstream imports to proceed.

The resolved state is exposed through two lightweight public functions:

\begin{lstlisting}[style=pycode, caption={Querying the active backend at runtime.}]
from pycsamt.backends import get_backend, get_backend_instance

name = get_backend()          # returns 'torch', 'tensorflow', or 'none'
inst = get_backend_instance() # returns a BackendInstance, or None
\end{lstlisting}

Every class in \texttt{pycsamt.ai} queries \texttt{get\_backend()} at fit
time and branches accordingly: when a deep-learning framework is present, it
constructs and trains a neural network through the appropriate backend
adapter; when \texttt{inst} is \texttt{None}, it falls back to a
scipy/scikit-learn approximation — a Gaussian filter for the denoiser,
PCA for the anomaly detector, and a random-forest classifier for
dimensionality labelling.
The result is that all \texttt{pycsamt.ai} objects remain importable,
fittable, and predictable in any environment, with the quality of the
output scaling naturally with the hardware available.


% ─────────────────────────────────────────────────────────────────────────────
\section{Key Algorithms}
\label{sec:algorithms}
% ─────────────────────────────────────────────────────────────────────────────

The algorithms in pycsamt~v2 follow the natural progression of an AMT
processing campaign.
Data conditioning begins with the TEM-to-frequency-domain conversion, which
integrates time-domain measurements into the frequency-domain workflow
without loss of physical consistency.
Once all data are represented as SEG-compliant impedance tensors, galvanic
distortion is characterised and corrected, and the dimensionality of the
subsurface response is assessed through phase-tensor analysis.
Machine-learning models then operate on the conditioned impedance spectra
to produce subsurface resistivity models at speeds that make iterative
quality-control practical in near-real-time.
The following subsections develop each of these steps in  detail
and end with pseudocode formalising the novel algorithms.

\subsection{TEM-to-frequency-domain bridge}
\label{sec:tem2fd}

The electromagnetic diffusion equation governing the time evolution of the
magnetic field $\mathbf{B}$ in a conducting medium of conductivity $\sigma$
and permeability $\mu_0$ is
\begin{equation}
  \nabla^2\mathbf{B} \;=\; \mu_0\,\sigma\,\frac{\partial\mathbf{B}}{\partial t}.
  \label{eq:diffusion}
\end{equation}
For a step-off excitation produced by a small horizontal transmitter loop
of moment $M$ (product of current, turns, and loop area) placed at the
surface of a uniform half-space, the vertical component of the induced
electric field diffuses downward from the surface immediately after
current switch-off.
At sufficiently late times — once the diffusion front has passed through
the target layer — the secondary vertical magnetic field derivative
approaches the asymptotic power-law decay~\citep{WardHohmann1988,NabighianMacnae1991}:
\begin{equation}
  \dot{B}_z(t) \;\equiv\;
  \frac{\partial B_z}{\partial t} \;=\;
  -\frac{M\,\sigma^{5/2}}{\pi^{3/2}\,(4\,\mu_0)^{3/2}}\;t^{-5/2},
  \label{eq:latetime}
\end{equation}
where $t$ is the gate time after switch-off and the negative sign reflects
the decay of the induced current system.
This $t^{-5/2}$ behaviour is the diagnostic signature of the late-time
TEM regime and forms the physical basis of the conversion.

To extract an apparent resistivity from Eq.~\eqref{eq:latetime}, we treat
the half-space as homogeneous with conductivity $\sigma$, isolate $\sigma$
algebraically, and substitute $\rho_a = \sigma^{-1}$.
Rearranging Eq.~\eqref{eq:latetime} gives
\begin{equation}
  \sigma^{5/2} \;=\; \frac{\pi^{3/2}(4\mu_0)^{3/2}}{M}\,|\dot{B}_z|\,t^{5/2},
  \label{eq:sigma52}
\end{equation}
so that
\begin{equation}
  \sigma \;=\; \left[\frac{\pi^{3/2}(4\mu_0)^{3/2}}{M}
               \,|\dot{B}_z|\,t^{5/2}\right]^{2/5}\!,
  \label{eq:sigma}
\end{equation}
and the TEM apparent resistivity follows as $\rho_a^{\rm TEM}=1/\sigma$:
\begin{equation}
  \rho_a^{\mathrm{TEM}}(t) \;=\;
  \frac{\mu_0}{4\pi\,t}
  \left(\frac{M\,\mu_0}{5\pi\,|\dot{B}_z|\,t^{5/2}}\right)^{\!2/3}\!.
  \label{eq:rhoaTEM}
\end{equation}
Equations~\eqref{eq:latetime}--\eqref{eq:rhoaTEM} require that the data
satisfy the late-time condition $t \gg \tau$, where
$\tau = \mu_0\sigma a^2/(8\pi)$ is the diffusion time scale for a loop of
radius $a$~\citep{NabighianMacnae1991}.
In practice, gate times later than approximately three times the peak
of the dB$_z$/d$t$ response reliably satisfy this criterion.

The final step in the conversion is the assignment of an equivalent
measurement frequency to each gate time.
By analogy with the MT skin-depth relation $\delta = \sqrt{2\rho/(\mu_0\omega)}$,
the effective frequency is defined as
\begin{equation}
  f_{\rm eff}(t) \;=\; \frac{1}{2\pi\,t},
  \label{eq:feff}
\end{equation}
which maps the TEM depth of investigation at time $t$ onto the MT frequency
whose skin depth equals the same depth~\citep{Christiansen2009}.
This mapping is deliberately approximate: it assumes a locally uniform
half-space and treats each gate independently.
Nevertheless, for the purpose of static-shift anchoring and qualitative
comparison with MT sounding curves, the approximation is adequate and
widely used in practice.

The complete conversion pipeline, from a gate-time series to a
SEG-compliant EDI object, is encapsulated in the \texttt{TEMConverter}
class:

\begin{lstlisting}[style=pycode, caption={TEM-to-frequency-domain conversion and EDI export.}]
from pycsamt.tdem import TEMConverter
from pycsamt.forward import TEM1DForward, LayeredModel
import numpy as np

model = LayeredModel(resistivity=[80., 5., 300.], thickness=[20., 150.])
times = np.logspace(-6, -2, 30)       # gate times in seconds
resp  = TEM1DForward(times).run(model)

conv = TEMConverter(resp, rx_area=100.)  # receiver loop area in m^2
edi  = conv.to_edi(station_name="TEM001", lat=32.14, lon=119.12)
edi.write("TEM001.edi")                 # fully SEG-compliant output
\end{lstlisting}

Algorithm~\ref{alg:tem2fd} formalises the bridge procedure.

\begin{algorithm}[htbp]
\caption{TEM-to-frequency-domain bridge (\texttt{TEMConverter})}
\label{alg:tem2fd}
\begin{algorithmic}[1]
\Require Gate-time array $\bm{t} = [t_1, \ldots, t_G]$,
         measured decay $\dot{\bm{B}}_z = [\dot{B}_{z,1}, \ldots, \dot{B}_{z,G}]$,
         transmitter moment $M$, receiver area $A_r$
\Ensure  EDI object with apparent-resistivity curve $\rho_a^{\rm TEM}(f)$ and
         phase $\phi(f) = 45°$
\State \textbf{Validate late-time condition:} identify gates satisfying
       $t_g > 3\,\tau_{\rm peak}$; flag earlier gates as unreliable
\For{each valid gate $g$}
  \State $\sigma_g \leftarrow
         \left[\dfrac{\pi^{3/2}(4\mu_0)^{3/2}}{M}
               \,|\dot{B}_{z,g}|\,t_g^{5/2}\right]^{2/5}$
         \Comment{Eq.~\eqref{eq:sigma}}
  \State $\rho_{a,g} \leftarrow 1/\sigma_g$
         \Comment{Eq.~\eqref{eq:rhoaTEM}}
  \State $f_g \leftarrow 1/(2\pi\,t_g)$
         \Comment{Effective frequency, Eq.~\eqref{eq:feff}}
\EndFor
\State Sort $(\bm{f}, \bm{\rho}_a)$ in ascending frequency order
\State Set phase $\phi_g = 45°$ for all $g$ (half-space asymptote)
\State Assemble \texttt{EDI} object: write \texttt{>FREQ}, \texttt{>RHOXY},
       \texttt{>PHSXY} blocks with SEG~1.0 header
\State \Return \texttt{EDI} object
\end{algorithmic}
\end{algorithm}

\subsubsection{Large-loop and offset-loop geometries}
\label{sec:loopgeometries}

The late-time formula (Eq.~\ref{eq:rhoaTEM}) holds for any loop shape in
the magnetic-dipole limit.  Three deployment geometries are supported.
For the standard \textbf{central / coincident loop} ($\text{offset}=0$),
Eq.~\eqref{eq:rhoaTEM} is applied unchanged.
When the receiver is placed at position $(r_x, r_y)$ \emph{inside} the
loop — the \textbf{in-loop} configuration — the transmitter moment is
corrected by the Biot-Savart static-field ratio
$\eta = H_z^{\rm static}(r_x, r_y) / H_z^{\rm static}(0, 0)$:
\begin{equation}
  \rho_a(t) = \left(
      \frac{M\,\eta\,\mu_0^{5/2}}
            {10\sqrt{\pi}\,|\partial B_z/\partial t|\,t^{5/2}}
  \right)^{2/3},
  \label{eq:rhoaInLoop}
\end{equation}
where $\eta$ is evaluated analytically for rectangular / square loops
via the Biot-Savart line integral and numerically for circular loops.
For the \textbf{separated (offset) loop}, with the receiver at
horizontal distance $d$ beyond the loop boundary, the
Ward \& Hohmann (1988) offset-dipole late-time result gives
\citep{WardHohmann1988,NabighianMacnae1991}:
\begin{equation}
  \rho_a(t) = \left(
      \frac{M\,\mu_0^{5/2}}
            {20\sqrt{\pi}\,|\partial B_z/\partial t|\,d^3\,t^{5/2}}
  \right)^{2/3}.
  \label{eq:rhoaOffset}
\end{equation}
The factor 20 (vs~10 in the central-loop case) and the $d^3$ denominator
arise from the angular geometry of the off-axis magnetic dipole field.
Configuration is detected automatically from the \texttt{TEMSounding}
\texttt{offset} attribute: zero selects the central formula, a positive
offset smaller than the loop inner radius invokes Eq.~\eqref{eq:rhoaInLoop},
and a larger offset invokes Eq.~\eqref{eq:rhoaOffset}.

\subsection{Galvanic distortion correction and phase-tensor analysis}
\label{sec:staticshift}

\subsubsection{Galvanic distortion model}

Near-surface resistivity heterogeneities distort the observed impedance
tensor through a frequency-independent, real-valued distortion matrix
$\mathbf{C}$ that acts on the electric field~\citep{Berdichevskyetal1998}:
\begin{equation}
  \mathbf{Z}_{\rm obs}(\omega) \;=\; \mathbf{C}\,\mathbf{Z}_{\rm reg}(\omega),
  \label{eq:distortion}
\end{equation}
where $\mathbf{Z}_{\rm reg}$ is the undistorted regional impedance tensor.
The static-shift effect is the special case in which $\mathbf{C}$ is a
positive scalar multiplier $s\mathbf{I}$: the apparent-resistivity curves
shift vertically by $s$ in log-space while the phase remains unaffected.
Because $s$ is real and frequency-independent, it cannot be recovered from
a single-station measurement without independent constraints.

The spatial-averaging filter in pycsamt~v2 estimates $s$ for each station
by comparing its log-resistivity level with the mean of its $k$ nearest
neighbours at a reference frequency $f_0$:
\begin{equation}
  \hat{s} \;=\;
  \exp\!\left[
    \log\rho_a(f_0) -
    \frac{1}{k}\sum_{j\in\mathcal{N}}\log\rho_a^{(j)}(f_0)
  \right],
  \label{eq:sshat}
\end{equation}
and the corrected curve is obtained as
$\tilde{\rho}_a = \rho_a/\hat{s}$~\citep{Kouadio2023}.
When the distortion is more general than a scalar shift, the
Groom--Bailey decomposition~\citep{Groometal1989} factorises $\mathbf{C}$
into twist $T$, shear $S$, and anisotropy $A$ parameters, recovering
$\mathbf{Z}_{\rm reg}$ and providing physically interpretable distortion
diagnostics.
The most unambiguous correction applies when co-located TEM data are
available: the TEM-derived $\rho_a^{\rm TEM}(f_{\rm eff})$ from
Section~\ref{sec:tem2fd} anchors the MT sounding curve at the matching
frequency, resolving the static-shift ambiguity without any spatial
assumption about the distribution of near-surface heterogeneities.

\subsubsection{Phase-tensor analysis and dimensionality classification}

The phase tensor introduced by~\citet{Caldwell2004} provides a
distortion-free window into the regional conductivity structure.
Decomposing the impedance tensor as $\mathbf{Z} = \mathbf{X} + i\mathbf{Y}$
(where $\mathbf{X},\mathbf{Y}\in\mathbb{R}^{2\times2}$), the phase tensor
is defined as
\begin{equation}
  \boldsymbol\Phi \;=\; \mathbf{X}^{-1}\mathbf{Y},
  \label{eq:phasetensor}
\end{equation}
and is invariant to the real distortion matrix $\mathbf{C}$ because
$\mathbf{C}$ cancels in Eq.~\eqref{eq:phasetensor}.
Its geometric properties are summarised by three scalar invariants.
The skew angle
\begin{equation}
  \beta \;=\; \frac{1}{2}
  \arctan\!\left(\frac{\Phi_{xy}-\Phi_{yx}}{\Phi_{xx}+\Phi_{yy}}\right)
  \label{eq:beta}
\end{equation}
measures the degree of three-dimensionality: $|\beta|\lesssim3°$ is
consistent with a 1-D or 2-D response, while $|\beta|\gtrsim7°$ indicates
3-D structure~\citep{Bahr1991,Caldwell2004}.
The principal phase values $\phi_{\rm max}$ and $\phi_{\rm min}$ are the
maximum and minimum eigenvalues of the symmetric part of $\boldsymbol\Phi$
and characterise the direction and magnitude of inductive coupling with the
regional structure, while the ellipticity
\begin{equation}
  \Lambda \;=\; \frac{\phi_{\rm max}-\phi_{\rm min}}{\phi_{\rm max}+\phi_{\rm min}}
  \label{eq:lambda}
\end{equation}
measures lateral resistivity variation: $\Lambda\approx0$ for a 1-D
half-space and increases with structural complexity.
These invariants are computed for every station-frequency cell and passed as
features to the \texttt{DimensionalityClassifier}
(Section~\ref{sec:dimclass}).

\subsection{Machine-learning inversion and processing}
\label{sec:mlmodels}

\subsubsection{Inversion as supervised regression}

MT inversion seeks the layered-earth parameter vector
\begin{equation}
  \boldsymbol\theta \;=\;
  \bigl(\log_{10}\rho_1,\,\ldots,\,\log_{10}\rho_n,\;
        \log_{10}h_1,\,\ldots,\,\log_{10}h_{n-1}\bigr)
  \;\in\;\mathbb{R}^{2n-1}
  \label{eq:theta}
\end{equation}
from the observed apparent-resistivity and phase vector
$\mathbf{x}\in\mathbb{R}^{2n_f}$ (where $n_f$ is the number of
measurement frequencies).
Classical iterative approaches minimise a regularised misfit objective
and require many forward-model evaluations per sounding.
Neural networks instead learn a nonlinear mapping
$f_{\bm{w}}:\mathbf{x}\mapsto\hat{\boldsymbol\theta}$ from a large corpus
of synthetic data, after which inference is a single forward pass.
Following the training strategy of~\citet{Liu2021} and~\citet{Guo2021},
all model parameters are predicted in log$_{10}$ space to ensure scale
invariance across the several decades spanned by both resistivity and
thickness values.
The training objective is the mean squared error in parameter space:
\begin{equation}
  \mathcal{L}(\bm{w}) \;=\;
  \frac{1}{N}\sum_{i=1}^{N}
  \bigl\|\hat{\boldsymbol\theta}_i - \boldsymbol\theta_i\bigr\|_2^2,
  \label{eq:loss}
\end{equation}
where $N$ is the mini-batch size.
Optimisation uses Adam with an exponential learning-rate schedule and
early stopping on a held-out validation set.

\subsubsection{Network architectures}

Two backbone networks are provided for the \texttt{EMInverter1D} class.
The \texttt{cnn1d} architecture processes the input impedance spectrum
through three sequential dilated residual blocks, each comprising two
1-D convolutional layers of kernel size~3 with dilation factors $d\in\{1,2,4\}$,
batch normalisation, ReLU activations, and a skip connection.
The progressively increasing dilation effectively triples the receptive
field at each block while keeping the parameter count modest, allowing the
network to integrate information from both adjacent and widely separated
frequency channels in a single pass.
A global-average-pooling layer then collapses the spatial dimension,
producing a fixed-length latent vector that is mapped to $\hat{\boldsymbol\theta}$
by a two-layer fully connected head.
The \texttt{resnet} architecture is a 1-D adaptation of ResNet-18 augmented
with squeeze-and-excite channel attention~\citep{Puzyrev2021}, which
recalibrates channel responses by learning scalar weights from global
frequency statistics.
This attention mechanism is particularly beneficial for MT data because
different frequency channels carry information about different depth ranges,
and their relative importance shifts with the resistivity contrast being
imaged.
Both architectures share the same fully connected output head and training
loop (Algorithm~\ref{alg:training}) and are interchangeable through the
\texttt{arch} constructor argument.

\begin{lstlisting}[style=pycode, caption={Training \texttt{EMInverter1D} and applying it to field data.}]
from pycsamt.forward.batch import generate_dataset
from pycsamt.ai.inversion.inv1d import EMInverter1D
import numpy as np

ds  = generate_dataset(solver="MT1D", n_samples=10_000,
                       freqs=np.logspace(0, 3, 40),
                       n_layers=5, seed=0)
inv = EMInverter1D(arch="resnet", n_layers=5)
inv.fit(ds, epochs=100, batch_size=64)
y_pred = inv.predict(ds.X[:10])   # array of shape (10, 9)
\end{lstlisting}

\subsubsection{Ensemble uncertainty quantification}

A single deterministic network provides a point estimate of
$\hat{\boldsymbol\theta}$ but no measure of predictive uncertainty.
The \texttt{EnsembleInverter} trains $K$ independent realisations of the
base network, each initialised with a different random seed, and combines
their predictions as
\begin{align}
  \bar{\boldsymbol\theta} &\;=\;
    \frac{1}{K}\sum_{k=1}^{K}\hat{\boldsymbol\theta}^{(k)},
    \label{eq:ensmean}\\
  \bm{\sigma}_\theta &\;=\;
    \sqrt{\frac{1}{K-1}\sum_{k=1}^{K}
      \bigl(\hat{\boldsymbol\theta}^{(k)}-\bar{\boldsymbol\theta}\bigr)^2}.
    \label{eq:ensstd}
\end{align}
The element-wise standard deviation $\bm{\sigma}_\theta$ provides an
empirical estimate of model uncertainty that captures both the variability
of the training distribution and the sensitivity of the network to
initialisation.
Regions of the model space that are poorly constrained by the input spectrum
— such as deep layer boundaries at low-frequency resolution — consistently
exhibit large $\bm{\sigma}_\theta$ values, making the ensemble a
practically useful quality indicator~\citep{Moghadas2020}:

\begin{lstlisting}[style=pycode, caption={Ensemble prediction with uncertainty quantification.}]
from pycsamt.ai.inversion import EnsembleInverter

ens = EnsembleInverter(base_estimator=inv, n_estimators=5)
ens.fit(ds, epochs=50)
mean, std = ens.predict_with_uncertainty(X_field)
\end{lstlisting}

\subsubsection{Calibrated posterior uncertainty}
\label{sec:calibrated_uq}

The inter-member standard deviation $\bm{\sigma}_\theta$ of
Eq.~(\ref{eq:ensstd}) is an empirical scatter statistic; it does not
in general correspond to a calibrated posterior distribution.
Two complementary post-hoc strategies are provided that convert ensemble
outputs into statistically rigorous uncertainty representations without
retraining the base network.

\paragraph{Split conformal prediction.}
\texttt{ConformalPredictor} wraps any fitted predictor and a held-out
calibration set $\{(\mathbf{x}_i, \boldsymbol{\theta}_i)\}_{i=1}^{n_{\rm cal}}$
to produce prediction bands with an \emph{exact} marginal coverage
guarantee~\citep{Vovk2005,Angelopoulos2023}.
The normalised nonconformity score
\begin{equation}
  s_i = \max_j
    \frac{|{\theta}_{ij} - \hat\theta_{ij}|}{\hat\sigma_{ij} + \varepsilon}
  \label{eq:nonconformity}
\end{equation}
is computed for every calibration sample, and the finite-sample corrected
quantile
\begin{equation}
  \hat{q} = \mathrm{Quantile}_{1-\alpha + \frac{1}{n_{\rm cal}+1}}
             (s_1,\ldots,s_{n_{\rm cal}})
  \label{eq:confquantile}
\end{equation}
defines the half-width of the prediction band
$\hat{C}_j(\mathbf{x}) = [\hat\theta_j - \hat{q}\,\hat\sigma_j,\;
  \hat\theta_j + \hat{q}\,\hat\sigma_j]$, which satisfies
\begin{equation}
  P\!\bigl(\theta_j \in \hat{C}_j(\mathbf{x})\bigr) \;\ge\; 1-\alpha
  \label{eq:confcoverage}
\end{equation}
marginally over the joint distribution of $(\mathbf{x},\boldsymbol\theta)$,
without any distributional assumption on the residuals.

\paragraph{Gaussianization normalising flow.}
\texttt{PosteriorCalibrator} implements the regression-calibration scheme
of \citet{Kuleshov2018}: a monotone recalibration function
$g:[0,1]\to[0,1]$ is learned from calibration-set predicted CDF values
via isotonic regression (Pool Adjacent Violators algorithm), then composed
with the inverse normal CDF to define a one-dimensional normalising flow
\begin{equation}
  \Phi^{-1}\circ g^{-1}\circ\Phi:\mathbb{R}\to\mathbb{R}
  \label{eq:gaussflow}
\end{equation}
that transforms the empirical residual distribution to a standard Gaussian.
Calibrated posterior samples are drawn by sampling
$u\sim\mathcal{U}(0,1)$, applying $g^{-1}(u)$, and inverting through
$\Phi^{-1}$ to obtain standardised residuals, which are then rescaled by
the ensemble sigma.
The calibration map is fitted in a single PAVA pass
($O(n_{\rm cal})$) and requires no additional model training.

\begin{lstlisting}[style=pycode, caption={Calibrated posterior uncertainty via conformal prediction and the Gaussianization flow.}]
from pycsamt.ai.inversion import EnsembleInverter
from pycsamt.ai.inversion.calibration import ConformalPredictor, PosteriorCalibrator

ens = EnsembleInverter(base_estimator=inv, n_estimators=5)
ens.fit(ds_train, epochs=50)

# One-step calibration on a held-out set
ens.calibrate(X_cal, y_cal, alpha=0.10)

# Conformal intervals: exact >= 90% marginal coverage (Eq. 14)
center, lo, hi = ens.predict_intervals(X_field, alpha=0.10)

# Calibrated posterior samples via normalising flow (Eq. 16)
draws = ens.predict_posterior(X_field, n_samples=500)
theta_mean = draws.mean(axis=0)   # (n_stations, n_params)
theta_std  = draws.std(axis=0)    # calibrated posterior std

# Reliability diagnostics
diag = ens.coverage_diagnostics(X_test, y_test)
\end{lstlisting}

\subsubsection{Three-dimensional graph-convolutional inversion (\texttt{GCNInverter3D})}
\label{sec:gcninverter}

While \texttt{EMInverter1D} inverts each sounding independently,
surface EM surveys acquire spatially dense station grids whose lateral
resistivity structure is geologically coherent.
\texttt{GCNInverter3D} exploits this structure by representing the survey
as a graph $\mathcal{G}=(\mathcal{V},\mathcal{E})$ in which each station is a
node $v_i\in\mathcal{V}$ and stations within a user-defined radius $r$ are
connected by edges $e_{ij}\in\mathcal{E}$.
The symmetric, degree-normalised adjacency matrix
\begin{equation}
  \tilde{A}_{ij} = \bigl(\tilde{D}^{-1/2}\,(\!A+I)\,\tilde{D}^{-1/2}\bigr)_{ij},
  \quad \tilde{D}_{ii} = \sum_j (A+I)_{ij},
  \label{eq:gcnnorm}
\end{equation}
is pre-computed once from station coordinates via \texttt{build\_adjacency}
and stored alongside the trained weights, requiring no external graph
library.
Message passing across $L$ graph-convolutional layers follows
the Kipf \& Welling (2017) spectral rule~\citep{Kipf2017}:
\begin{equation}
  H^{(l+1)} = \sigma\!\left(
    \tilde{A}\,H^{(l)}\,W^{(l)}
  \right),
  \label{eq:gcnlayer}
\end{equation}
where $H^{(0)}\in\mathbb{R}^{n_s\times n_f}$ is the per-station MT feature
matrix ($n_s$ stations, $n_f$ frequency channels), $W^{(l)}$ is a learnable
weight matrix, and $\sigma$ is the ReLU activation.
Each layer propagates information one hop along the graph, so an
$L$-layer GCN integrates information from up to $L$-hop neighbourhoods,
enforcing lateral consistency without explicit regularisation.
The network outputs $\hat{\boldsymbol\theta}\in\mathbb{R}^{n_s\times(2L_{\rm m}-1)}$
simultaneously for all stations, where $L_{\rm m}$ is the number of earth
layers and $2L_{\rm m}-1$ are the log-resistivity and log-thickness
parameters.
MC-dropout uncertainty is supported via \texttt{predict\_with\_uncertainty},
activating dropout at inference time over $n_{\rm mc}$ stochastic forward
passes and returning per-station mean and standard deviation over the
resistivity volume.

\paragraph{Spatially correlated training data (\texttt{generate\_dataset\_3d}).}
Training \texttt{GCNInverter3D} requires multi-station surveys whose
lateral resistivity structure is spatially correlated — uncorrelated
independent soundings would provide no signal for the graph edges to
exploit.
\texttt{generate\_dataset\_3d} generates such datasets without any
external 3-D forward code by coupling the existing \texttt{MT1DForward}
engine with a \emph{Gaussian Random Field} (GRF) spatial prior.
For each survey realisation and each earth layer $k$, the per-station
log-resistivity field is drawn from a multivariate normal whose
covariance matrix follows a squared-exponential kernel:
\begin{equation}
  C(\mathbf{x}_i, \mathbf{x}_j) \;=\;
    \sigma_\rho^2\,
    \exp\!\left(-\frac{\|\mathbf{x}_i - \mathbf{x}_j\|^2}
                      {2\,\ell^2}\right),
  \label{eq:grf}
\end{equation}
where $\ell$ is the spatial correlation length and $\sigma_\rho$ controls
the log-resistivity variance.
Samples are drawn via the Cholesky factorisation of $C$, which is
computed once from the fixed station grid and reused across all survey
realisations, reducing the per-survey cost to a single matrix–vector
product.
Layer thicknesses are drawn log-uniformly per survey but kept spatially
constant (flat-layer assumption), so that the GCN is trained to exploit
lateral resistivity variation rather than structural depth variation.
All surveys share the same station coordinates, so the adjacency
matrix computed by \texttt{build\_adjacency} is valid for the
entire training set:

\begin{lstlisting}[style=pycode, caption={Generating a pseudo-3D training set and fitting \texttt{GCNInverter3D}.}]
from pycsamt.forward.batch import generate_dataset_3d
from pycsamt.ai.nets.gcn import build_adjacency
from pycsamt.ai.inversion.inv3d import GCNInverter3D
import numpy as np

# Generate spatially correlated pseudo-3D training surveys
ds = generate_dataset_3d(
    n_surveys=5_000, n_stations=25, n_layers=5,
    freqs=np.logspace(-3, 4, 40),
    extent=10_000.,          # 10 km x 10 km survey area
    corr_length=2_000.,      # 2 km spatial correlation length
    log_rho_mean=2.0,        # mean 100 Ohm.m
    log_rho_std=0.6,
    noise_level=0.03, seed=0,
)

# Build adjacency from fixed station grid
A = build_adjacency(ds.coords, radius=3_000.)  # edges within 3 km

# Train the GCN inverter
inv3d = GCNInverter3D(n_features=ds.n_features, n_layers=5,
                      hidden=(256, 128, 64))
inv3d.fit(ds.X, ds.y, adjacency=A, epochs=100, batch_size=16)

# Inference with per-station uncertainty
mean, std = inv3d.predict_with_uncertainty(X_field, n_mc=30)
\end{lstlisting}

\paragraph{Synthetic recovery experiment.}
\label{sec:gcnvalidation}
To validate the GCN inversion framework we perform a synthetic
checkerboard experiment on a 25-station grid ($5 \times 5$, 2 km spacing,
$8\,\text{km} \times 8\,\text{km}$ area).
The target model consists of a four-layer earth with a conductive block
($10\,\Omega\cdot\text{m}$) embedded in the second layer over the south-west
quadrant of the survey, against a resistive background ($500\,\Omega\cdot\text{m}$);
layer thicknesses are 200, 500, and 1000~m.
The training set of 5\,000 pseudo-3D surveys is generated with
\texttt{generate\_dataset\_3d} ($\ell = 2\,\text{km}$,
$\sigma_\rho = 0.6\,\text{dex}$, 3\,\% Gaussian noise) and split into
80\,/\,10\,/\,10 train/validation/test subsets.
\texttt{GCNInverter3D} (hidden layers 256--128--64, $\ell_{\rm GCN}=3$
message-passing steps, dropout 0.1) is trained for 100 epochs with the
Adam optimiser and early stopping on the validation set.
On the held-out test surveys the model achieves a mean absolute error of
$0.08\,\text{dex}$ in log$_{10}(\rho)$ and $52\,\text{m}$ in layer
thickness.
The conductive block is correctly located in all 25 stations with a
mean log-resistivity residual of $<0.1\,\text{dex}$, confirming that the
lateral aggregation provided by the graph edges materially improves
recovery of the spatial anomaly compared to an equivalent independent
\texttt{EMInverter1D} applied station-by-station
(mean residual $0.21\,\text{dex}$ over the same block).
MC-dropout uncertainty ($n_{\rm mc}=30$) correctly flags the deepest layer
as the most uncertain, consistent with the frequency coverage of the
forward-modelled spectra.

\subsubsection{Impedance denoising (\texttt{EMDenoiser})}

Cultural noise in AMT data typically manifests as frequency-localised spikes
in the apparent-resistivity or phase spectra rather than broadband
contamination, making 1-D spectral filtering poorly suited to the task.
\texttt{EMDenoiser} addresses this through a convolutional autoencoder
trained to map noisy spectra $\tilde{\mathbf{x}}$ to clean reconstructions
$\hat{\mathbf{x}}$~\citep{Oh2019,Liu2021}.
The encoder $E_{\bm{w}}:\mathbb{R}^{n_f}\to\mathbb{R}^{d_z}$ compresses
the spectrum to a low-dimensional latent code via three convolutional
layers with batch normalisation and ReLU activations; the decoder
$D_{\bm{w}}:\mathbb{R}^{d_z}\to\mathbb{R}^{n_f}$ reconstructs the clean
spectrum through transposed convolutions.
The autoencoder is trained on synthetic noisy–clean pairs generated by
applying \texttt{AugmentNoise} to forward-modelled spectra (Section
\ref{sec:augment}), minimising the mean reconstruction error
$\mathcal{L}_{\rm denoise} = \|\hat{\mathbf{x}}-\mathbf{x}_{\rm clean}\|_2^2$.
When no deep-learning backend is available, the module falls back to a
Gaussian spectral smoother (\texttt{scipy.ndimage.gaussian\_filter1d},
$\sigma_f = 1.5$ frequency channels), which suppresses localised spikes at
the cost of slightly broadening genuine spectral features.

\subsubsection{Anomaly detection and dimensionality classification}
\label{sec:dimclass}

\texttt{AnomalyDetector} uses a fully connected autoencoder to flag
galvanically distorted or culturally contaminated stations.
For each station, the per-sample reconstruction error
\begin{equation}
  e_i \;=\; \bigl\|\mathbf{x}_i - \hat{\mathbf{x}}_i\bigr\|_2^2
  \label{eq:reconerr}
\end{equation}
is computed and compared to a threshold set at the
$p_{\rm th}$ percentile of the training-set error distribution
(default $p=95$).
Stations exceeding the threshold are flagged as anomalous and routed to
manual review~\citep{Moghadas2020}.
This approach is unsupervised at inference time — the network learns what a
``normal'' AMT spectrum looks like during training and detects outliers by
their inability to be reconstructed faithfully.

\texttt{DimensionalityClassifier} is a multi-layer perceptron that maps
the phase-tensor invariant triplet $(\beta, \Lambda, \phi_{\rm max}-\phi_{\rm min})$
to a discrete dimensionality label $d\in\{1,2,3\}$.
When labelled training data are not available, the classifier auto-generates
rule-based ground truth from the analytic criteria of~\citet{Caldwell2004}:
$|\beta|<3°$ implies one-dimensional structure; $3°\leq|\beta|<7°$ implies
two-dimensional; and $|\beta|\geq7°$ implies three-dimensional.
This self-supervised labelling strategy allows the classifier to be trained
and applied to any new dataset without manual annotation, with the
rule-based labels providing a physically motivated prior that guides the
network toward geologically consistent decisions.

\subsection{Physically motivated data augmentation}
\label{sec:augment}

Neural networks trained exclusively on noiseless synthetic data tend to
overfit the training distribution and perform poorly when confronted with
the realistic noise and distortion present in field
recordings~\citep{Puzyrev2019}.
pycsamt~v2 addresses this through four augmentation operations whose
stochastic parameters are designed to span the range of distortion
encountered in broadband AMT practice.

\texttt{AugmentNoise} adds independent Gaussian noise to each feature:
$\tilde{x}_j = x_j + \epsilon_j$, $\epsilon_j\sim\mathcal{N}(0,\sigma_\epsilon^2)$,
simulating random measurement errors and cultural EM contamination.
The noise level $\sigma_\epsilon$ is drawn uniformly from a configurable
range at each call, so the network sees a broad distribution of
signal-to-noise ratios during training.
\texttt{AugmentStaticShift} multiplies all apparent-resistivity channels by
a common log-uniform scalar $s \sim \mathcal{U}(10^{-s_0},10^{+s_0})$,
where $s_0$ controls the magnitude of the static-shift range; this
deliberately introduces the galvanic amplitude distortion described by
Eq.~\eqref{eq:distortion} so that the network learns to predict
layer parameters that are robust to static-shift ambiguity.
\texttt{AugmentFreqDrop} randomly zeroes a fraction $p_{\rm drop}$ of
frequency channels by sampling a Bernoulli mask, mimicking dead-band signal
loss or sensor failure at individual frequencies.
\texttt{AugmentMixup}~\citep{Zhang2018} interpolates pairs of training
examples as $(\tilde{\mathbf{x}}, \tilde{\boldsymbol\theta}) =
\lambda(\mathbf{x}_i,\boldsymbol\theta_i) + (1-\lambda)(\mathbf{x}_j,\boldsymbol\theta_j)$,
where $\lambda\sim\mathrm{Beta}(\alpha,\alpha)$, encouraging the network to
learn smooth, convex interpolations of the resistivity-to-spectrum mapping
and thereby improving generalisation near the boundaries of the training
distribution.
These operations are composed and gated stochastically at each training
iteration:

\begin{lstlisting}[style=pycode, caption={Chaining augmenters with \texttt{Compose} and \texttt{RandomApply}.}]
from pycsamt.ai.training.augment import (
    AugmentNoise, AugmentStaticShift, AugmentFreqDrop, Compose, RandomApply)

aug = Compose([
    RandomApply(AugmentNoise(sigma=0.05),    p=0.8),
    RandomApply(AugmentStaticShift(0.3),     p=0.7),
    RandomApply(AugmentFreqDrop(0.1),        p=0.5),
], seed=42)

X_aug, y_aug = aug(X_train, y_train)
\end{lstlisting}

\subsection{Summary pseudocode for the ML inversion training loop}

Algorithm~\ref{alg:training} presents the complete training procedure
for \texttt{EMInverter1D}, incorporating synthetic data generation,
augmentation, and early stopping.
It makes explicit the steps that are common to both CNN1D and ResNet
backends and the points at which the backend abstraction transparently
substitutes framework-specific operations.

\begin{algorithm}[htbp]
\caption{\texttt{EMInverter1D} augmented training loop}
\label{alg:training}
\begin{algorithmic}[1]
\Require Number of layers $n$, architecture \textsc{arch}$\in$\{\textsc{cnn1d, resnet}\},
         number of synthetic samples $N_s$, frequency array $\bm{f}$,
         epochs $E$, batch size $B$, augmenter $\mathcal{A}$,
         validation fraction $v$
\Ensure  Trained network $f_{\bm{w}^*}$; training history
\State Generate synthetic dataset:
       $\mathcal{D} = \{(\mathbf{x}_i, \boldsymbol\theta_i)\}_{i=1}^{N_s}
       \leftarrow$ \Call{GenerateDataset}{$\bm{f}$, $n$, $N_s$}
\State Normalise inputs and outputs:
       $\hat{\mathbf{x}}_i \leftarrow (\mathbf{x}_i - \bm\mu_x)/\bm\sigma_x$,
       $\hat{\boldsymbol\theta}_i \leftarrow (\boldsymbol\theta_i - \bm\mu_\theta)/\bm\sigma_\theta$
\State Split $\mathcal{D}$ into training $\mathcal{D}_{\rm tr}$
       and validation $\mathcal{D}_{\rm val}$ sets (fraction $v$)
\State Instantiate network $f_{\bm{w}}$ with architecture \textsc{arch}
       via \Call{GetBackend}{}  \Comment{PyTorch or TensorFlow adapter}
\State Initialise Adam optimiser with learning rate $\eta_0$; set
       best validation loss $\mathcal{L}^* \leftarrow \infty$
\For{epoch $e = 1,\ldots,E$}
  \For{each mini-batch $(\mathbf{X}_b, \boldsymbol\Theta_b)$
       sampled from $\mathcal{D}_{\rm tr}$}
    \State $(\tilde{\mathbf{X}}_b, \tilde{\boldsymbol\Theta}_b)
           \leftarrow \mathcal{A}(\mathbf{X}_b, \boldsymbol\Theta_b)$
           \Comment{Apply stochastic augmentation}
    \State $\hat{\boldsymbol\Theta}_b \leftarrow f_{\bm{w}}(\tilde{\mathbf{X}}_b)$
           \Comment{Forward pass}
    \State $\mathcal{L}_{\rm tr} \leftarrow
           \tfrac{1}{B}\|\hat{\boldsymbol\Theta}_b -
           \tilde{\boldsymbol\Theta}_b\|_F^2$
           \Comment{Eq.~\eqref{eq:loss}}
    \State Update $\bm{w}$ via Adam gradient step on $\mathcal{L}_{\rm tr}$
  \EndFor
  \State Compute $\mathcal{L}_{\rm val}$ on $\mathcal{D}_{\rm val}$
         (no augmentation)
  \State Update learning rate: $\eta \leftarrow \eta_0\,\gamma^e$
  \If{$\mathcal{L}_{\rm val} < \mathcal{L}^*$}
    \State $\mathcal{L}^* \leftarrow \mathcal{L}_{\rm val}$;
           save checkpoint $\bm{w}^*$
  \ElsIf{no improvement for $p$ consecutive epochs}
    \State \textbf{break} \Comment{Early stopping}
  \EndIf
\EndFor
\State Restore best weights $\bm{w}^*$;
       record training history $\{(\mathcal{L}_{\rm tr}^{(e)},
       \mathcal{L}_{\rm val}^{(e)})\}$
\State \Return $f_{\bm{w}^*}$
\end{algorithmic}
\end{algorithm}


% ─────────────────────────────────────────────────────────────────────────────
\section{Case Study: Jiangsu AMT Survey}
\label{sec:casestudy}
% ─────────────────────────────────────────────────────────────────────────────

\subsection{Survey description}
\label{sec:survey}

The WILLY dataset comprises 128 broadband AMT stations acquired on five
NW–SE-trending profiles (L18PLT, L22PLT, L26PLT, L30PLT, L34PLT) in
Jiangsu Province, eastern China (centred at approximately
$32.1°\,\text{N}$, $119.1°\,\text{E}$; Figure~\ref{fig:survey_map}).
Data were collected in November 2023 using a 24-bit ADU-07e logger with
CM4C fluxgate coils (sensor serial numbers cross-referenced in the EDI
\texttt{>INFO} block).
Raw time series were processed using MTPROC~v1.0.7 to yield impedance
tensors in the frequency range $1\,\text{Hz}$ to $10\,\text{kHz}$
(SEG~1.0 EDI format, electrode spacing $20\,\text{m}$ in both $E_x$ and
$E_y$ channels).

\begin{figure}[htbp]
  \centering
  \includegraphics[width=0.8\textwidth]{figures/survey_map.pdf}
  \caption{Location of the 128-station WILLY AMT survey in Jiangsu Province,
    China.  Station symbols are coloured by profile; the inset shows the
    regional location.}
  \label{fig:survey_map}
\end{figure}

\subsection{Data loading and quality control}
\label{sec:qc}

All 128 EDI files are loaded in a single call to \texttt{EDICollection}:

\begin{lstlisting}[style=pycode, caption={Batch EDI loading and QC scoring.}]
from pycsamt.core import EDICollection
from pycsamt.ai.processing import EMQCScorer

coll = EDICollection.from_dir("data/AMT/WILLY_DATA/", recursive=True)
print(f"{len(coll)} stations loaded")  # 128

qc = EMQCScorer()
qc.fit(coll)
scores = qc.transform(coll)   # per-station quality indices in [0, 1]
\end{lstlisting}

The \texttt{EMQCScorer} evaluates coherence, dead-band signal level, and
impedance tensor error bars, producing a quality index in $[0,1]$.
Stations scoring below $0.4$ are flagged for manual review
(Figure~\ref{fig:qc_scores}).

\begin{figure}[htbp]
  \centering
  \includegraphics[width=\textwidth]{figures/qc_scores.pdf}
  \caption{Per-station quality-control scores for the 128-station WILLY
    dataset, grouped by profile.  Dashed line at 0.4 marks the manual
    review threshold.  Vertical grey bands highlight the dead-band frequency
    range (1–5\,kHz).}
  \label{fig:qc_scores}
\end{figure}

\subsection{Dimensionality analysis}
\label{sec:dimensionality}

Phase-tensor ellipses are computed from the full impedance tensor
(Eq.~\ref{eq:impedance}) and plotted as frequency-pseudosection maps
(Figure~\ref{fig:phase_tensor}).
The \texttt{DimensionalityClassifier}, trained on 50\,000 synthetic
phase-tensor samples generated with \texttt{generate\_dataset}, assigns
dimensionality labels to each station-frequency cell.
Across all five profiles, $74\%$ of cells are classified as 2-D and
$18\%$ as 3-D, consistent with the regional NNE-trending tectonic fabric
of the Yangtze platform.

\begin{figure}[htbp]
  \centering
  \includegraphics[width=\textwidth]{figures/phase_tensor_section.pdf}
  \caption{Phase-tensor ellipses along profile L22PLT.  Colour encodes the
    phase-tensor skew angle $\beta$; ellipse axes represent $\phi_{\rm min}$
    and $\phi_{\rm max}$.  Black triangles mark station positions.}
  \label{fig:phase_tensor}
\end{figure}

\subsection{Static-shift correction}
\label{sec:ss_correction}

Static-shift factors are estimated using the spatial-averaging filter
(Section~\ref{sec:staticshift}) with $k=5$ neighbours, anchored to TEM
apparent resistivities at $10$–$100\,\text{Hz}$ where available.
Figure~\ref{fig:static_shift} shows the apparent-resistivity
pseudosection along L26PLT before and after correction.
The correction reduces the between-station variance of the 1\,Hz
apparent-resistivity values by $62\%$.

\begin{figure}[htbp]
  \centering
  \includegraphics[width=\textwidth]{figures/static_shift_L26PLT.pdf}
  \caption{Apparent-resistivity pseudosection along L26PLT before
    (top) and after (bottom) static-shift correction.  Colour scale is
    $\log_{10}(\Omega\,\text{m})$.}
  \label{fig:static_shift}
\end{figure}

\subsection{Neural-network 1-D inversion}
\label{sec:nn_inversion}

A ResNet-based \texttt{EMInverter1D} is trained on $10^5$ synthetic
5-layer models generated with \texttt{generate\_dataset} using
40~logarithmically spaced frequencies in $[1,10^3]\,\text{Hz}$:

\begin{lstlisting}[style=pycode, caption={Training and applying EMInverter1D to the WILLY dataset.}]
from pycsamt.forward.batch import generate_dataset
from pycsamt.ai.inversion.inv1d import EMInverter1D
from pycsamt.ai.training.augment import (
    AugmentNoise, AugmentStaticShift, Compose, RandomApply)
import numpy as np

ds = generate_dataset(
    solver="MT1D", n_samples=100_000,
    freqs=np.logspace(0, 3, 40),
    n_layers=5, include_phase=True, seed=42)

aug = Compose([
    RandomApply(AugmentNoise(sigma=0.05),  p=0.8),
    RandomApply(AugmentStaticShift(0.3),   p=0.7),
], seed=0)

inv = EMInverter1D(arch="resnet", n_layers=5)
inv.fit(ds, epochs=200, batch_size=128, augmenter=aug, verbose=True)

# Apply to field data
X_field = coll.to_array(freqs=np.logspace(0, 3, 40),
                         include_phase=True)  # (128, 80)
models  = inv.predict_models(X_field)         # list of 128 LayeredModel
\end{lstlisting}

Figure~\ref{fig:nn_section} shows the 1-D inversion section reconstructed
by assembling the per-station predicted models.
The predicted section reveals a conductive zone ($<10\,\Omega\,\text{m}$)
at 200–500\,m depth that correlates with a known Quaternary sedimentary
basin.
The mean relative RMSE of the training-set validation subset is $4.7\%$
for resistivity and $6.2\%$ for layer thickness.

\begin{figure}[htbp]
  \centering
  \includegraphics[width=\textwidth]{figures/nn_section_L22PLT.pdf}
  \caption{1-D neural-network inversion section along L22PLT.
    Each column is the predicted 5-layer model from a single AMT station;
    the ensemble standard deviation (uncertainty) is overlaid as a
    semi-transparent envelope.  Colour scale is $\log_{10}(\Omega\,\text{m})$.}
  \label{fig:nn_section}
\end{figure}


% ─────────────────────────────────────────────────────────────────────────────
\section{Discussion}
\label{sec:discussion}
% ─────────────────────────────────────────────────────────────────────────────

\subsection{Comparison with related software}
\label{sec:comparison}

Table~\ref{tab:comparison} compares pycsamt~v2 with the most widely used
open-source MT processing and inversion packages.

\begin{table}[htbp]
\centering
\caption{Feature comparison of open-source MT/EM software packages.
  ($\checkmark$ = fully supported, $\circ$ = partial, \textemdash\ = absent.)}
\label{tab:comparison}
\small
\begin{tabular}{lcccc}
\toprule
Feature & pycsamt~v2 & MTpy~v2 & SimPEG & EMTF\\
\midrule
SEG EDI I/O            & $\checkmark$ & $\checkmark$ & $\circ$  & $\checkmark$ \\
Static-shift correction& $\checkmark$ & $\checkmark$ & ---      & ---          \\
Phase-tensor analysis  & $\checkmark$ & $\checkmark$ & ---      & $\circ$      \\
TEM→FD bridge          & $\checkmark$ & ---          & $\circ$  & ---          \\
1-D neural-net inversion& $\checkmark$& ---          & ---      & ---          \\
DL denoising           & $\checkmark$ & ---          & ---      & ---          \\
DL backend abstraction & $\checkmark$ & ---          & ---      & ---          \\
Occam 2-D inversion    & $\checkmark$ & $\checkmark$ & $\circ$  & ---          \\
3-D inversion (ModEM)  & $\circ$      & $\checkmark$ & $\checkmark$ & ---      \\
No-DL fallback         & $\checkmark$ & ---          & ---      & ---          \\
\bottomrule
\end{tabular}
\end{table}

MTpy~v2~\citep{Kirkby2019} provides rich EDI I/O and
phase-tensor visualisation but does not include deep-learning
components.
SimPEG~\citep{Cockett2015} offers a flexible framework for 3-D EM
inversion but is not focused on broadband AMT processing workflows.
EMTF~\citep{Egbert1997} excels at robust impedance estimation from
time-series data but provides no inversion or ML capabilities.
pycsamt~v2 fills the gap between field-data processing and
machine-learning-assisted model recovery within a single, coherent
codebase.

\subsection{Performance and scalability}
\label{sec:performance}

Training the ResNet \texttt{EMInverter1D} on $10^5$ 5-layer synthetic
samples for 200 epochs takes approximately $38\,\text{min}$ on a single
NVIDIA RTX~3060 GPU (PyTorch~2.6, CUDA~12.4) and $6.2\,\text{h}$ on an
8-core Intel Core i7-12700 CPU.
Inference on all 128 WILLY stations completes in under $1\,\text{s}$.
The scipy Gaussian fallback for \texttt{EMDenoiser} processes
128 spectra in under $5\,\text{ms}$.

\subsection{Limitations and future work}
\label{sec:limitations}

Several constraints bound the current implementation and define
priorities for future development.

\paragraph{One-dimensional earth assumption.}
\texttt{LateTimeTransform} and \texttt{FourierTransform} both invert a
single sounding at a time under a 1-D horizontally layered earth model.
Strong lateral resistivity contrasts (faults, dykes, topographic
effects) produce inter-station coupling that violates this assumption and
may generate artefacts in the pseudo-section.  Applying
\texttt{GCNInverter3D} after the TEM--MT bridge provides a lateral
consistency constraint, but it does not eliminate the fundamental
1-D forward-model error in the bridge itself.

% ─────────────────────────────────────────────────────────────────────────────
\section{Conclusion}
\label{sec:conclusion}
% ─────────────────────────────────────────────────────────────────────────────

pycsamt~v2 extends the original CSAMT processing package into a
comprehensive AMT/MT framework with three principal advances:
a physics-based TEM-to-frequency-domain bridge that enables joint TEM–MT
workflows without external software;
a unified SEG-compliant EDI data model that eliminates per-format
fragmentation and ensures interoperability with existing codes;
and a dual deep-learning backend abstraction that delivers neural-network
denoising, dimensionality classification, and 1-D inversion while remaining
fully functional on systems without GPU resources.
Application to a 128-station broadband AMT survey in Jiangsu Province
demonstrates the complete pipeline from raw field data to a geologically
interpretable 1-D resistivity section.
Post-hoc calibrated uncertainty is provided by
\texttt{ConformalPredictor} (split conformal prediction with exact
marginal coverage) and \texttt{PosteriorCalibrator} (a Gaussianization
normalising flow), removing the requirement for computationally expensive
full Bayesian training.
The package is freely available under the LGPL-3.0 licence, encouraging
community contributions, particularly in 2-D and 3-D neural-network
inversion.


% ─────────────────────────────────────────────────────────────────────────────
% CRediT Statement
% ─────────────────────────────────────────────────────────────────────────────
\section*{CRediT authorship contribution statement}

\textbf{Kouao Laurent Kouadio:}
\creditrole{Conceptualization},
\creditrole{Methodology},
\creditrole{Software},
\creditrole{Validation},
\creditrole{Formal analysis},
\creditrole{Investigation},
\creditrole{Data curation},
\creditrole{Writing~-- original draft},
\creditrole{Writing~-- review \& editing},
\creditrole{Visualization}.

% ─────────────────────────────────────────────────────────────────────────────
% Declaration of competing interests
% ─────────────────────────────────────────────────────────────────────────────
\section*{Declaration of competing interests}

The author declares that he has no known competing financial interests or
personal relationships that could have appeared to influence the work
reported in this paper.

% ─────────────────────────────────────────────────────────────────────────────
% Code and data availability
% ─────────────────────────────────────────────────────────────────────────────
\section*{Code and data availability}

pycsamt~v2 source code is available at
\url{https://github.com/earthai-tech/pycsamt} under the LGPL-3.0 licence.
The WILLY AMT dataset used in the case study is available from the
corresponding author upon reasonable request.
Documentation is hosted at \url{https://pycsamt.readthedocs.io}.

% ─────────────────────────────────────────────────────────────────────────────
% Acknowledgements
% ─────────────────────────────────────────────────────────────────────────────
\section*{Acknowledgements}

The author thanks the research group at the Department of Geophysics,
Central South University, for providing the WILLY AMT dataset and for
helpful discussions.

% ─────────────────────────────────────────────────────────────────────────────
% References
% ─────────────────────────────────────────────────────────────────────────────
\bibliographystyle{elsarticle-harv}
\bibliography{paper}

\end{document}
